\section{Discussion}
\label{sec:discussion}

Table~\ref{tab:hyp} answers the four hypotheses the project committed to before any results.

\begin{table}[t]
\centering
\caption{Pre-registered hypotheses and outcomes (SD\,1.5 unless stated).}
\label{tab:hyp}
\footnotesize
\setlength{\tabcolsep}{3pt}
\begin{tabular}{p{0.9cm}p{3.0cm}p{3.6cm}}
\toprule
H & Claim & Outcome\\
\midrule
H1 & $\leq$50 pairs suffice to generalise to the five held-out slides &
\textbf{Supported for SD\,1.5} with source conditioning and DDIM inversion (every slide improves);
not for SDXL at 1024\,px, which needed 144 pairs.\\
\midrule
H2 & Colour improves over raw input and classical baselines, on MITOS and CAMELYON17 &
\textbf{Partly.} Beats raw everywhere and classical on typical slides; loses to Macenko on A06 and to
ParamNet/StainGAN everywhere; CAMELYON17 fails for every SD\,1.5 configuration, final pipeline
included, and passes narrowly for SDXL.\\
\midrule
H3 & Relative Dice $\geq 0.95$ &
\textbf{Supported with fusion}: 0.994 for inversion + F3; not without it (0.893 unfused, 0.875
for the A4-type point).\\
\midrule
H4 & 4-step LCM matches 50-step DDIM &
\textbf{Not supported}: generic LCM is 12.5$\times$ faster but loses quality; task-specific distillation
is worse.\\
\bottomrule
\end{tabular}
\end{table}

\textbf{Why the learned baselines win.} ParamNet and StainNet apply a colour map to the input pixels, and
StainGAN's CycleGAN generator keeps a wide feature map (a quarter of the resolution, hundreds of
channels); none of them passes the image through a 48$\times$ compressing four-channel latent. Every diffusion output is resynthesised by the VAE, which alone limits SSIM to
0.576--0.646 on Aperio crops (Section~\ref{sec:vae}). F3 fusion is, in effect, our way of turning the
diffusion model back into a colour map applied to the source: its success on structure and the
learned colour maps' success on colour point in the same direction. With 50 pairs, the strong image prior
of a diffusion model does not compensate for the information it loses in its own bottleneck.

\textbf{Colour metrics.} LAB Wasserstein compares colour histograms and ignores layout, so histogram
matching wins it on A06 by construction while scoring worse than doing nothing on typical slides. The
windowed variant and per-slide reporting were essential to avoid reading that as a real gain. Pixelwise
$\Delta E_{2000}$ tracks SSIM almost rank for rank, so it is not an independent colour check.

\textbf{Cost.} Timed on one RTX~3090 with loading excluded (Table~\ref{tab:app-speed}), the final pipeline
takes 6.27\,s per 512\,px crop, generic LCM at four steps 0.50\,s (12.5$\times$ faster), and SDXL 12.1\,s per
1024\,px crop, about 3.0\,s per 512\,px of tissue. The learned baselines take 2--25\,ms and the classical
ones 6--42\,ms. A 100{,}000$\times$100{,}000 pixel whole-slide image holds about 38{,}000 crops of
512\,px, so the final pipeline needs about 66 GPU-hours per slide, LCM 5.3 and SDXL 32, against
minutes for every baseline; the few-step LCM path that was meant to remove this cost did not hold its
quality (H4).

\textbf{Limitations.} All training data comes from one slide pair (A03/H03), so slide-level variation is
not represented in training. The learned baselines were trained once, with one seed each. Structural
safety rests on one segmenter (HoVer-Net) and one colon dataset (Lizard). The planned clinical-utility test has no valid
instrument: every atypia classifier we built failed its own pre-registered gate
(Appendix~\ref{app:clinical}), so no clinical-utility claim is made. The decoder smoothing analysis
rests on two inspected crops.

\textbf{Planned but not run.} Several planned items were not carried out and no result is claimed for
them: the mitosis-detection arm of the clinical-utility evaluation, which needs a separate object
detector; the SD\,3.5 probe's structural-ControlNet and SD\,3.5 Turbo few-step checks; a full-VAE fine-tune
with the LoRA and ControlNet retrained on its new latents (only the anchored version was run), and a cross-centre check of the fine-tuned decoder on CAMELYON17; and a
source-image conditioning branch for PixCell. Two choices also departed from the plan: the TCGA-BRCA colour
reference was not used as a secondary colour target, and the
attention-residual method~\cite{madusanka2025attentionresidual}, planned as the state-of-the-art colour
baseline, was not reproduced, so it is cited but not compared.
